Дуже легко втратити час і енергію, коли духовний пошук скерований на речі, що мають мало або й жодного значення для пробудження. Три провідні ідеї дають понятійний каркас, на якому спочиває вчення, і скеровують розум до ключових принципів, що ведуть до просвітлення. Ці провідні ідеї дають фокус і напрям трьом основним практикам, які описано далі в цій книзі.

\section*{Питання буття}

Над входом до Дельфійського оракула було написано слова «Пізнай самого себе». Прийшов Ісус і додав до цієї давньої ідеї відчуття нагальності й наслідків, коли сказав: «Якщо ти винесеш назовні те, що в тобі, те, що ти винесеш, урятує тебе. Якщо ти не винесеш назовні те, що в тобі, те, що ти не винесеш, знищить тебе».

Ісус каже, що духовність — річ серйозна й із серйозними наслідками. Ваше життя хистко тримається рівноваги, хитаючись між станами несвідомого сновидіння наяву та повної уважності. Те, що більшість людей не сприймають своє життя так, засвідчує, наскільки глибоко вони ще сплять і заперечують.

То що ж це те, що нам потрібно винести назовні?

Усередині кожної з наших форм спочиває екзистенційна таємниця \emph{буття}. За фізичною подобою, особистістю, статтю, історією, заняттям, надіями й мріями, приходами й відходами лежить таємнича тиша, безодня спокою, сповнена ефірної присутності. Навіть із усім нашим тривожним клопотом і одержимістю дрібницями ми не можемо цілком зігнорувати цю примарну сутність у своєму центрі. І все ж ми робимо все, що можемо, щоб уникнути її спокою, її тиші, її порожнечі та її сяйної інтимності.

\emph{Буття} — це те, що заважає нашій наполегливості залишатися в царстві заціпеніння нашого потаємного відчаю. Це свербіж, який не можна почухати, свист, який не можна ігнорувати. Буття, справжнє буття — це не просто вірування.

Більшість із нас живе в стані, у якому наше \emph{буття} давно заслане в похмуре царство мовчазної туги. Часом \emph{буття} розриває тканину нашої несвідомості, щоб нагадати нам, що ми не живемо тим життям, яким могли б жити, — життям, яке справді має значення. Іншими митями \emph{буття} відступає за лаштунки, мовчки чекаючи нашої відданої уваги. Але не обманюйтеся: \emph{буття} — ваше \emph{буття} — це осередкова тема життя.

Залишатися несвідомим щодо \emph{буття} — це бути замкненим у спустошеній і сповненій конфліктів землі, керованій его, яка здається нормальною, бо ми дозволили завести себе в стан незнання, де товстий шар ненависті, нечесності й жадібності вважається нормальним і здоровим. Але насправді ніщо не могло б бути менш здоровим, ніж те, що ми, люди, називали дійсністю.

Тримаючись того, що ми знаємо й у що віримо, ми прив'язуємося до руху своєї уяви й зумовленого мислення, водночас вважаючи себе цілком раціональними й здоровими. Тож ми далі зміцнюємо дійсність того, що завдає стільки болю й страждання нам самим та іншим.

Глибоко в душі ми всі підозрюємо, що щось дуже не так із тим, як ми сприймаємо життя, але ми робимо все, щоб утекти від цього факту. І так само — те, як ми залишаємося сліпими до свого жахливого стану через одержиме й патологічне заперечення \emph{буття} — наче з нами станеться щось жахливе, якщо ми глянемо в чисте світло Істини й поховаємо свої прив'язаності до ілюзії.

Саме у вимірі \emph{буття} відкривається Істина — не істина математики чи хімії, філософії чи історії, а Істина, яка починає розгортатися в ті миті спокою, коли звична рутина життя раптом стає прозорою для відчуття сенсу й значущості, невідомого більшість часу. Ці життєво важливі й несподівані зустрічі з \emph{буттям} вказують на Істину, що лежить одразу під тканиною нашого повсякденного життя, нагадуючи нам, що життя, за яке ми тримаємося, може бути дурнішим, ніж ми будь-коли уявляли, і що є Дійсність, яка має силу відімкнути таємницю наших життів, якщо ми вирішимо відмовитися від своєї відданості страхові й безпеці, які так добре знаємо.

Усі ми народжуємося з \emph{буттям}, закритим темрявою. Ми можемо впізнати прозорість \emph{буття}, що сяє в очах немовляти, але таке \emph{буття} не усвідомлює себе. Немовлята живуть у чарівному світі несвідомого \emph{буття}, тоді як дорослі живуть у світі егоцентричної відокремленості й заперечення \emph{буття}. Духовне пробудження уможливлює відновлення й повернення \emph{буття} до його справжньої влади.

Питання \emph{буття} — це все. Ніщо не могло б бути важливішим чи наслідковішим — ніщо, де ставки такі високі. Залишатися несвідомим щодо \emph{буття} — це залишатися заціпенілим щодо власної дійсності, а отже, заціпенілим щодо Дійсності в широкому масштабі. Вибір простий: прокинутися до \emph{буття} або спати сном без кінця.
